%Indique ac\'a, las herramientas de LLMs/IA usadas en su proyecto. \\En caso de no utilizar ninguna explicitar: \textbf{Declaro no haber utilizado ninguna herramienta de LLMs/IA en el desarrollo de este trabajo.}\\ \noindent En caso de s\'i utilizarlo, para cada herramienta/uso, indicar:
\begin{enumerate} 
\item \textbf{Herramienta utilizada:} \textit{Gemini Pro} y \textit{NotebookLM} en sus versiones gratuitas para estudiantes.
\item \textbf{Propósito de uso:} 

\begin{itemize}
\item Generar las bibliografías del archivo \texttt{.bib}.
\item Revisión de redacción y ortografía. 
\item Generar listados de comparación entre modelos según la literatura. 
\item Escritura en \LaTeX{} de fórmulas matemáticas.
\item Implementación de scripts para la generación de gráficas.
\end{itemize}

\item \textbf{Prompts relevantes:} 

\begin{itemize}
\item ``Genera el formato BibTeX de la referencia [X]''.
\item ``Compara el modelo de la fuente [X] con el de la fuente [Y]''.
\item ``Dado el siguiente archivo en \LaTeX{}, indica errores ortográficos y de redacción en tercera persona''.
\item ``Escribe en formato \LaTeX{} la ecuación [i] del artículo [X]''.
\item ``En \texttt{code/algorithms/plots.py}, implementa una función que genere gráficas de [x] vs. [y] con formato monocromático''.
\end{itemize}

\item \textbf{Nivel de edición:} Reescritura y revisión directa en el documento para la comparación de modelos. Incorporación manual de entradas \texttt{.bib} y fórmulas en \LaTeX{}. Las sugerencias de redacción se evaluaron y seleccionaron manualmente.

\item \textbf{Validación:} Lectura y contraste directo con las fuentes originales (\textit{NotebookLM} indica la sección específica de consulta). Revisión manual y ejecución del código en \textit{Python} para los gráficos generados.

\end{enumerate}


%Herramienta: ChatGPT (GPT-5.x)
%Uso: Generación de ideas y revisión de redacción
%Prompt principal: “Explica el algoritmo SMS-EMOA con pseudocódigo formal”
%Intervención: El texto fue adaptado, reorganizado y complementado con material de clase
%Validación: Se contrastó con [fuente X]
